% GDZ PPN236018647,native573middle;HKUST574.
\NumberedItem{100}
\indent Circa terminos medios arithmetico-geometricos multa nova deteximus.
\par\noindent\hfill Br[unsvigae, 1799] Novemb.

\medskip\begingroup\small
\indent Man wird als Belege für die hier angezeigten Entdeckungen die Leisteaufzeichnungen im Abschnitt [II.], oben S. 177 und die artt. [1.]—[3.] aus der Scheda Ac des Abschnitts [IV.], oben S. 184 anzusehen haben. Natürlich sind nicht alle diese Entdeckungen im November gemacht worden, die Eintragung ins \textit{Tagebuch} wird vielmehr nur die Bedeutung haben, daß vor der Abreise nach Helmstedt ein Teil der erlangten Ergebnisse in die damals begonnene Scheda Ac eingetragen worden ist.
\par\noindent\hfill\textsc{Klein.\quad Schlesinger.}
\par\endgroup

\NumberedItem{101}
\indent Medium arithmetico-geometricum tamquam quotientem duarum functionum transscendentium repraesentabile esse iam pridem inveneramus; nunc alteram harum functionum ad quantitates integrales reducibilem esse deteximus.
\par\noindent\hfill Helmst[adii, 1799] Dec. 14.

\medskip\begingroup\small
\indent Der erste Satz bezieht sich auf die Quotientendarstellung [7], oben S. 186, der zweite auf die Formeln [13]—[16] des art. [4.], oben S. 187, nach denen der Zähler dem reziproken Werte von
\[
\int_0^1\frac{dr}{\sqrt{(1-r^2)(x^2-r^2)}}
\]
\noindent gleich gefunden wird. Das ursprünglich irrtümlicherweise hingeschriebene Integral zweiter Gattung hat \textsc{Gauss} vergessen durchzustreichen.
\par\noindent\hfill\textsc{Klein.\quad Schlesinger.}
\par\endgroup
